\section{Conclusion and Future Work}
\label{sec:conclusion}

\paragraph{Achieved results.}
We presented a new proof rule for almost-sure termination of probabilistic while-loops with loop-free bodies, phrased entirely in distribution-transformer semantics. The rule separates two certificates: a non-negative supermartingale that bounds the probability of escaping to ever-larger sublevel sets, and a natural-valued variant that need only decrease with a uniformly positive probability and be bounded within each such sublevel set. We proved the rule sound and relatively complete by relating program-level certificates to the pCFG-level rule of Majumdar and Sathiyanarayana~\cite{DBLP:journals/pacmpl/MajumdarS25}: soundness follows by compiling our certificates into pCFG certificates, and relative completeness follows in the converse direction. On the one-dimensional random walk, the rule reproduces the textbook AST argument without any location-indexed bookkeeping. For the Fast Dice Roller, we gave AST proofs both at pCFG level and at program level; the comparison shows that the program-level proof is markedly shorter, reuses an existing distributional invariant, and avoids introducing artificial constants purely to enforce descent across intermediate control-flow edges. Finally, we proved almost-sure termination of the Fast Loaded Dice Roller, whose rejection-and-restart structure is captured by a single compact iteration-level argument: the supermartingale bounds the depth of the underlying decision tree, while the variant tracks the distance towards a non-rejecting leaf and resets, but does not grow unboundedly, on rejection. Combined with the existing partial-correctness result for such samplers~\cite{zilken2026verifyingsamplingalgorithmsdistributional}, our AST proofs complete the total-correctness argument for both FDR and FLDR. In particular, the FLDR proof closes the formal termination gap left by the original algorithm presentation.

\paragraph{Limitations.}
The rule as presented has two main restrictions. First, it applies to a single while-loop with a \emph{loop-free} body: the one-iteration distribution-transformer step that the rule reasons about is only well defined when the body itself contains no further loops. Nested loops can in principle be handled by applying the rule to the innermost loop first and treating its established AST as a black box when analyzing the enclosing loop, but we do not formalize or evaluate this modular composition here. Second, the underlying programming language does not include \emph{nondeterministic} (demonic or angelic) choice: every probabilistic branch has a fixed, syntactically given probability, and there is no construct for resolving nondeterminism via a scheduler. This mirrors the pCFG rule's treatment of nondeterminism-free models and keeps the connection between program syntax and one-step distribution-transformer semantics direct, but it excludes programs whose termination behavior depends on an adversarial or best-case choice of successor. As with the pCFG rule of~\cite{DBLP:journals/pacmpl/MajumdarS25}, we also restrict attention to programs over rational-valued variables, which avoids measure-theoretic subtleties but excludes continuous sampling primitives.

\paragraph{Future work.}
Lifting the loop-free-body restriction to a fully compositional treatment of nested and sequentially composed loops is the most immediate extension, and would let the rule be applied automatically along a program's syntax tree rather than loop by loop. Incorporating demonic or angelic nondeterminism, following strategy-based semantics for probabilistic programs with nondeterministic choice~\cite{DBLP:journals/pacmpl/BatzBKW24}, would connect the rule to settings where a scheduler resolves nondeterministic branching, as already supported by the pCFG rule of Majumdar and Sathiyanarayana~\cite{DBLP:journals/pacmpl/MajumdarS25}. It would also be worthwhile to study the precise relationship between our rule and the program-level rule of McIver et al.~\cite{10.1145/3158121} in more depth, in particular whether one can characterize classes of programs for which one of the two styles of certificate is systematically easier to construct than the other. Finally, the relative-completeness results now available for both pCFG-level and program-level AST rules~\cite{DBLP:journals/pacmpl/MajumdarS25,10.1145/3158121} suggest a concrete target for automation: since our rule decouples the supermartingale and variant obligations, each may be amenable to separate, specialized synthesis procedures, offering a path towards (semi-)automated AST proofs for sampling algorithms such as FDR and FLDR.
